\section{Massive tightness and the kernel count}
\label{sec:counting}

The preceding section accounts for all norm at bounded physical radii.
We first show that bounded massive graph norms cannot escape to infinity
as \(h\to\infty\). We then compare the resulting profiles with the exact
massive index. Throughout this section the induction hypotheses at sizes
less than \(N\) remain in force, and the exterior estimate at size \(N\)
has already been proved. The dimension of \(\cK_N\) is still unknown.

Write
\[
 \mathfrak q_{h,m}(\psi)
      =\frac1{16}\sum_v\norm{D_v(h,m)\psi}^2
\]
for the form of \(\cH(h,m)\), where the sum is over an orthonormal real
spin basis.

\subsection{Uniform confinement on large annuli}

\begin{samepage}
\begin{proposition}[Massive tightness]\label{prop:massive-tightness}
There are \(a>0\), \(h_0>0\), and \(R_0<\infty\), depending on \(N\),
such that
\begin{equation}\label{count:annular-confinement}
 \mathfrak q_{h,1}(\psi)\ge aR^2\norm{\psi}^2
 \quad\text{if}\quad
 h\ge h_0,\quad R\ge R_0,\quad
 \supp\psi\subset\{R<|x|<4R\}.
\end{equation}
This inequality is on gauge invariants, with no rotation restriction.
Consequently, for \(R\ge R_0\), after enlarging the constants if needed,
\begin{equation}\label{count:tail-bound}
 \int_{|x|>4R}|\psi|^2
 \le\frac{C}{R^2}\mathfrak q_{h,1}(\psi)
             +\frac{C}{R^4}\norm{\psi}^2,
 \qquad h\ge h_0.
\end{equation}
In particular, every bounded sequence of gauge- and \(G\)-invariant
vectors with \(h\to\infty\) and bounded \(D_u(h,1)\) graph norms has tight
squared norm measures.
\end{proposition}
\end{samepage}

Averaging the abelian charge squares cancels their angular terms. In the
application, each spinor charge acts on a different fast projection of the
same state. The following lemma shows that norm convergence of these
projections suffices for the cancellation.

\begin{lemma}[Center confinement for nearby states]
\label{count:center-confinement}
Fix a center space \(\mathcal Z_C\), an arbitrary spectator Hilbert space
\(\mathcal E\), and a compact set
\(K\subset\{c\in\mathcal Z_C:|c|\ge b_0\}\), where \(b_0>0\).
Let \((v_\nu)_{\nu=1}^{16}\) be an orthonormal real spin basis and
\(m_j>0\) with \(m_j\to\infty\). For each \(\nu\), let
\[
 a_{\nu,j}\in L^2(\mathcal Z_C;\mathcal F_C\otimes\mathcal E)
\]
be supported in \(K\), bounded in norm, and in the graph domain of the
center charge \(D_{v_\nu}^C(m_j)\), acting trivially on \(\mathcal E\).
Suppose that
\[
 \norm{a_{\nu,j}-a_{1,j}}\longrightarrow0,
 \qquad
 \norm{D_{v_\nu}^C(m_j)a_{\nu,j}}=o(m_j)
 \quad(\nu=1,\ldots,16).
\]
Then \(\norm{a_{\nu,j}}\to0\) for every \(\nu\).
No gauge or rotation invariance is imposed on these center states.
\end{lemma}

\begin{proof}
For brevity write \(a_{v,j}\) for the state belonging to a spinor in the
chosen basis, and put \(a_{0,j}=a_{v_1,j}\). The center mass multiplier
is bounded by \(Cm_j\) on \(K\). Subtracting it from the charge and using
the free center Dirac norm identity gives
\begin{equation}\label{count:center-gradient-bound}
 \norm{p_Ca_{v,j}}=\sqrt2\norm{D_v^C(0)a_{v,j}}=O(m_j).
\end{equation}
This identity holds with spectators by integration, or first on finite
spectator sums and then by closure.

Apply the abelian square formula of Proposition~\ref{prop:square} and
divide by \(m_j^2\). Dropping its nonnegative kinetic term yields
\begin{equation}\label{count:abelian-mass-lower-bound}
 0\ge \frac12\sum_{i=1}^9s_i^2\norm{c_i a_{v,j}}^2
       +\ip{a_{v,j}}{A_{v,j}a_{v,j}}+o(1),
 \qquad A_{v,j}=m_j^{-1}p_C\cdot R(v)c.
\end{equation}
The fermion multiplier and the spin part of the rotation generator
contribute only \(O(m_j^{-1})\). The vector field \(R(v)c\) has zero
divergence, so \(A_{v,j}\) is symmetric on these compactly supported
\(H^1\) states.

Bounded center coordinates and \eqref{count:center-gradient-bound} give
\[
 \norm{A_{v,j}a_{v,j}}+\norm{A_{v,j}a_{0,j}}\le C.
\]
In the second term, the free Dirac identity for \(v_1\) bounds the full
center gradient of \(a_{0,j}\), and hence bounds the angular operator
for every \(v\). Symmetry therefore gives
\begin{equation}\label{count:angular-replacement}
 \begin{split}
 &\left|\ip{a_{v,j}}{A_{v,j}a_{v,j}}
              -\ip{a_{0,j}}{A_{v,j}a_{0,j}}\right|\\
 &\hspace{1cm}\le\norm{a_{v,j}-a_{0,j}}
       \bigl(\norm{A_{v,j}a_{v,j}}+\norm{A_{v,j}a_{0,j}}\bigr)=o(1).
 \end{split}
\end{equation}
Only norm convergence of the difference is needed here. Bounded
coordinate multiplication also permits replacement by \(a_{0,j}\) in
the mass term with error \(o(1)\).

Average \eqref{count:abelian-mass-lower-bound} over the spin basis.
The identity \(\sum_\nu R(v_\nu)=0\) cancels the common angular
expectation exactly. Since \(s_i^2\ge1\),
\[
 0\ge\tfrac12\norm{|c|a_{0,j}}^2+o(1)
   \ge\tfrac12b_0^2\norm{a_{0,j}}^2+o(1).
\]
Thus \(a_{0,j}\to0\) in norm. The assumed norm differences give the
same conclusion for every \(a_{v,j}\).
\end{proof}

\begin{proof}[Proof of Proposition~\ref{prop:massive-tightness}]
Suppose that the eventual uniform assertion
\eqref{count:annular-confinement} is false. Its quantified negation gives
\(h_j\ge j\), \(R_j\ge j\), and unit vectors supported in
\(R_j<|x|<4R_j\) such that
\(\mathfrak q_{h_j,1}(\psi_j)/R_j^2\to0\). Dilation \(x=R_j\xi\),
performed unitarily, and multiplication of the charge by \(R_j\) give
parameters
\begin{equation}\label{count:annular-parameters}
 h'_j=h_jR_j^3,\qquad m'_j=R_j^2,\qquad
 \frac{h'_j}{m'_j}=h_jR_j\longrightarrow\infty.
\end{equation}
Writing the dilated unit states as \(f_j\), we have
\[
 \supp f_j\subset\{1<|\xi|<4\},\qquad
 \mathfrak q_{h'_j,m'_j}(f_j)=o((m'_j)^2).
\]
The sixteen nonnegative summands in this form imply
\(\norm{D_v(h'_j,m'_j)f_j}=o(m'_j)\) for each member of the spin basis.

First suppose, along a subsequence, that
\(m'_j/\sqrt{h'_j}\to\infty\). The exact scaling of the bosonic potential
is
\begin{equation}\label{count:potential-scaling}
 U_{h',m'}(\xi)=\frac{(m')^4}{(h')^2}
                    U_{1,1}\left(\frac{h'}{m'}\xi\right).
\end{equation}
Lemma~\ref{lem:massive-potential} and \eqref{count:annular-parameters}
therefore imply
\(U_{h'_j,m'_j}(\xi)\ge c(m'_j)^2|\xi|^2\) on the fixed annulus for all
large \(j\). The fermionic multiplier in Proposition~\ref{prop:square}
has operator norm at most \(C(h'_j+m'_j)\) there, and
\[
 \frac{h'_j+m'_j}{(m'_j)^2}\longrightarrow0.
\]
Dropping the nonnegative kinetic energy contradicts the assumed
\(o((m'_j)^2)\) form bound.

Otherwise pass to a subsequence with \(m'_j\le C\sqrt{h'_j}\). This is
the growing-mass regime of Lemmas~\ref{lem:raw-estimates} and
\ref{lem:fast-projection}. The norm measures concentrate on the commuting
part of the compact annulus. Every such configuration has at least two
joint centers. Cover this compact commuting set by finitely many
separated-center charts, chosen small enough for the fast estimates for
all sixteen spinors simultaneously. At each reference configuration its
block-center vector is nonzero. Shrinking the corresponding chart gives
\begin{equation}\label{count:center-away-zero}
                     |c|\ge b_0>0
\end{equation}
on it, with a possibly different positive constant in each of the finite
charts. A fixed smooth gauge-invariant partition localizes to these
charts and to a complementary region disjoint from the commuting set.
Its charge commutators are \(O(1)=o(m'_j)\). The complementary norm tends
to zero by the raw commutator estimate.

Consider one localized sequence, denoted by \(\psi_j\). In its full
\((b,t)\) Hilbert space set
\(a_{v,j}=P_v(b)\psi_j\). Lemma~\ref{lem:fast-projection} gives
\begin{equation}\label{count:variable-projections}
 \norm{a_{v,j}-\psi_j}=o(1),\qquad
 \norm{D_v^s(h'_j,m'_j)a_{v,j}}=o(m'_j).
\end{equation}
We retain the fast Hilbert space as a spectator here. The slow center
operator and the massive internal block operators anticommute, since
they act on disjoint variables and odd Clifford factors. Their squared
norm identity, first on a smooth compact slow core and then by closure,
implies
\begin{equation}\label{count:center-small-charge}
                    \norm{D_v^C(m'_j)a_{v,j}}=o(m'_j).
\end{equation}
This does not assert that \(P_v\) commutes with each block operator.
Rather, it applies the identity for the slow operator to the vector
\(a_{v,j}\), with its full dependence on \(b,t\). Compact slow support
and local ellipticity ensure the individual slow domains, just as in
\eqref{prof:block-graph-identity}.

The projections act only in the fast variables, so all the
\(a_{v,j}\) are supported in the same compact center set as the localized
states. By \eqref{count:center-away-zero}, this set is bounded away
from zero. Regard the internal variables and the fast Hilbert space as
spectators. Equations~\eqref{count:variable-projections} and
\eqref{count:center-small-charge} give every hypothesis of
Lemma~\ref{count:center-confinement}, with mass \(m'_j\).
The lemma therefore makes every \(a_{v,j}\), and hence \(\psi_j\),
tend to zero in norm. Summing the finitely many charts and using the
vanishing complementary norm contradicts \(\norm{f_j}=1\).
This proves \eqref{count:annular-confinement} in both parameter regimes.

For the tail estimate, use a smooth squared radial partition consisting
of an interior cutoff and annular functions \(\chi_j\), supported in
\(R_j<r<4R_j\) with \(R_j\ge R\) and geometrically increasing radii.
It can be chosen with bounded overlap, with the exterior functions
exhausting the partition for \(r>4R\), and
\(\sum_j|\nabla\chi_j|^2+|\nabla\chi_{\mathrm{in}}|^2\le C/R^2\).
Averaging the unweighted Dirac localization identity over the spin basis
gives
\[
 \mathfrak q_{h,1}(\chi_{\mathrm{in}}\psi)+
       \sum_j\mathfrak q_{h,1}(\chi_j\psi)
 \le\mathfrak q_{h,1}(\psi)+CR^{-2}\norm{\psi}^2.
\]
The interior form is nonnegative, and the exterior forms are bounded
below by \(aR^2\sum_j\norm{\chi_j\psi}^2\). This proves
\eqref{count:tail-bound}. Graph approximation and monotone exhaustion of
the partition justify the same argument for arbitrary form-domain
vectors. Finally, Lemma~\ref{lem:coercive-charge} gives
\(\mathfrak q_{h,1}(\psi)\le\tfrac32\norm{D_u(h,1)\psi}^2\) on gauge
and \(G\) invariants. The displayed tail bound proves their tightness.
\end{proof}

\subsection{Parity and the exact massive kernel dimension}

\begin{lemma}[Gap at large coupling]\label{lem:large-coupling-gap}
On the gauge- and \(G\)-invariant Hilbert space, there are \(\delta>0\)
and \(h_1<\infty\) such that, for all \(h\ge h_1\),
\begin{equation}\label{count:odd-gap}
 \norm{D_u(h,1)\psi}\ge\delta\norm{\psi}
                         \qquad(\psi\text{ odd}).
\end{equation}
If \(P_h\) is the orthogonal projection onto \(\Ker D_u(h,1)\) on that
space, then
\begin{equation}\label{count:full-gap}
 \norm{(1-P_h)\psi}\le\delta^{-1}\norm{D_u(h,1)\psi}.
\end{equation}
Moreover, the massive kernel is entirely even and has dimension \(p(N)\).
\end{lemma}

\begin{proof}
Let \(h_j\to\infty\), \(\norm{\psi_j}=1\), and
\(\norm{D_u(h_j,1)\psi_j}\to0\), with the prescribed invariances.
Proposition~\ref{prop:massive-tightness} and
Corollary~\ref{prof:full-exhaustion} give norm-exhausting profiles. By
Proposition~\ref{prop:profile-liminf}, every proper profile satisfies
\(D_u^C(1)F_C=0\) in the full center graph domain.

Each internal kernel factor is a rotation singlet by
Proposition~\ref{prop:singlet-parity}. The phase and symmetry of the fast
line in Lemma~\ref{lem:fast-line} and the compatibility in
Lemma~\ref{lem:chart-compatibility} therefore make \(G\) act on these
profiles through its ordinary action on the center variables and center
fermions. The abelian instance of Lemma~\ref{lem:coercive-charge}, valid
on the full graph domain by closure, gives
\[
       0=\norm{D_u^C(1)F_C}^2
          \ge\tfrac23\ip{F_C}{\cH_C(1)F_C}.
\]
This is understood as a form inequality. The mass-only oscillator has
a unique zero-energy line for each center color, as also follows from
the zero-triple case of Lemma~\ref{lem:normal-oscillators}. In orthonormal
center coordinates its bosonic factor is proportional to
\[
        \exp\left(-\frac12\sum_{i=1}^9|s_i|\,|c_i|^2\right),
\]
and its fermionic factor is the empty \(T\)-polarized vacuum. Denote the
normalized product by \(\Omega_C\). Its parity is even. Hence each proper
profile belongs to
\begin{equation}\label{count:proper-profile-kernel}
 \left(\C\Omega_C\otimes\bigotimes_{\alpha=1}^k\cK_{n_\alpha}\right)
                  ^{\text{equal-size permutations}}.
\end{equation}
The full-collapse profile belongs to \(\cK_N\) and is also even, by
Proposition~\ref{prop:singlet-parity}, regardless of the dimension of that
space.

If every \(\psi_j\) were odd, symmetry transport would make every profile
odd. The preceding identifications would force them all to vanish,
contradicting \eqref{prof:full-norm-exhaustion}. The negation of
\eqref{count:odd-gap} would give exactly such a sequence, by choosing
\(h_j\ge j\) and an odd unit vector with charge norm below \(1/j\).
Thus the eventual uniform odd gap holds.

For each fixed positive \(h\), Proposition~\ref{prop:massive-index} gives
self-adjointness and compact graph embedding on this invariant space.
The square has a discrete spectrum with finite-dimensional eigenspaces.
On the eigenspace of any positive eigenvalue \(\lambda\) of the square,
\(D_u/\sqrt\lambda\) is an isometric bijection between its even and odd
parts. The odd lower bound therefore excludes \(0<\lambda<\delta^2\)
in both parities. Spectral expansion proves
\eqref{count:full-gap}. There are no odd zero vectors; the index formula
\eqref{eq:massive-index} consequently gives
\(\dim\Ker D_u(h,1)=p(N)\), all even, for \(h\ge h_1\).
\end{proof}

\subsection{Partition multiplicities and recovery sequences}

The proper strata are indexed by the unordered partitions of \(N\) with
at least two parts. To see why there is only one stratum for each such
partition, order its block sizes in a fixed way and attach a distinct
center to each listed size. This labeled presentation is identified
only under permutations of entries of equal size. Relabeling unequal
sizes changes the presentation of the same commuting configuration; it
does not define a new partition type. Thus there are exactly
\(p(N)-1\) proper types.

By the induction hypothesis, each internal kernel in
\eqref{count:proper-profile-kernel} has dimension one. Choose the same
normalized kernel vector whenever the block sizes agree. Those vectors
are even. Exchanging two such factors consequently produces no graded
sign. The center Gaussian is unchanged by the corresponding orthogonal
permutation, and the empty center vacuum is fixed by its natural
exterior action. Lemma~\ref{lem:fast-line} supplies the same untwisted
action for the fast section, with positive volume normalizations.
Therefore \eqref{count:proper-profile-kernel} is exactly one dimensional
for each proper partition. Let its normalized vector be \(\Phi_C\).

\begin{lemma}[Recovery of the proper profile lines]
\label{count:proper-recovery}
There are even gauge- and \(G\)-invariant vectors \(v_{C,h}\), one for
each proper partition, such that, as \(h\to\infty\),
\begin{equation}\label{count:proper-recovery-estimates}
 \norm{D_u(h,1)v_{C,h}}\longrightarrow0,
 \qquad
 \ip{v_{C,h}}{v_{C',h}}\longrightarrow\mathbf1_{C=C'}.
\end{equation}
These vectors can be chosen simultaneously with recovery vectors for
any prescribed finite list in \(\cK_N\), with all cross inner products
tending to zero.
\end{lemma}

\begin{proof}
We first approximate \(\Phi_C\) in the center graph norm by admissible
tests. Multiply its center factor by a smooth outer radial cutoff
\(\zeta_R\), equal to one for \(|c|\le R\), zero for \(|c|\ge2R\), and
with derivative bounded by \(C/R\). Multiply also by the collision cutoff
\(\rho_\varepsilon\) of Lemma~\ref{lem:collision-removal}, chosen as a
product over all unordered pairs. These cutoffs are rotation invariant
and have the required permutation symmetry. The resulting vector
\(f_{C,R,\varepsilon}\) belongs to \(\mathcal T_C\).

Since \(D_u^C(1)\Phi_C=0\), only cutoff commutators contribute to its
charge. The outer error tends to zero by the Gaussian tail. For fixed
\(R\), the smooth center factor is bounded on its support, so
\eqref{prof:collision-capacity} bounds the squared collision error by
\(C_R\varepsilon^7\). Its omitted \(L^2\) norm also tends to zero, since
the collisions have measure zero. Consequently there is a sequence of
such tests \(f_{C,\ell}\) with
\begin{equation}\label{count:center-test-approximation}
 \norm{f_{C,\ell}-\Phi_C}\longrightarrow0,
 \qquad \norm{D_u^C(1)f_{C,\ell}}\longrightarrow0.
\end{equation}
The order is explicit: for tolerance \(1/\ell\), choose the outer radius
first and then choose the collision radius small enough for that fixed
outer support.

For each fixed \(\ell\), Proposition~\ref{prop:test-lifts} produces
\(w_{C,h,\ell}=I_{h,u,1}f_{C,\ell}\), with the prescribed leading norm
and charge output. We impose the remaining symmetries carefully, since
the chartwise corrections need not have been chosen rotation invariant.
Let \(A_G\) be Haar averaging over \(G\), and let
\(E_+=(1+(-1)^F)/2\) be the even projection. For every fixed \(g\in G\),
Lemma~\ref{lem:chart-compatibility} and invariance of the profile give
\[
 \ip{w_{C,h,\ell}}{U_gw_{C,h,\ell}}
                          \longrightarrow\norm{f_{C,\ell}}^2.
\]
These inner products are bounded uniformly in \(g,h\) for fixed
\(\ell\). Haar dominated convergence and the projection identity
\(\norm{A_Gw}^2=\ip{w}{A_Gw}\) show
\(\norm{w_{C,h,\ell}-A_Gw_{C,h,\ell}}\to0\). The identical argument
for the two-element parity action gives
\(\norm{w_{C,h,\ell}-E_+A_Gw_{C,h,\ell}}\to0\).
Moreover, \(D_u\) commutes with \(A_G\) and interchanges the parity
projections, so
\[
 \norm{D_u(h,1)E_+A_Gw_{C,h,\ell}}
                  \le\norm{D_u(h,1)w_{C,h,\ell}}.
\]
Thus imposing these symmetries preserves the leading norms and does not
increase the charge errors.

For fixed \(\ell\), the compact center supports on distinct partition
strata are disjoint compact subsets of the gauge quotient. The leading
identifications concentrate on those supports, so
Lemma~\ref{lem:chart-compatibility} makes the different lifts
asymptotically orthogonal. Their projections have the same property by
the norm approximation just proved.

There are only finitely many proper partitions. Choose an increasing
sequence \(H_\ell\to\infty\) so that for every \(h\ge H_\ell\) the finitely
many norm, overlap, symmetry, and graph-output errors at the fixed test
level \(\ell\) are below \(1/\ell\). Taking that level for
\(H_\ell\le h<H_{\ell+1}\), and normalizing the vectors, proves
\eqref{count:proper-recovery-estimates}. Additional finitely many cross
overlaps can be included when choosing \(H_\ell\); the required
whole-block vectors and their concentration are constructed next.
\end{proof}

\begin{lemma}[Recovery of any finite whole-block list]
\label{count:whole-recovery}
Let \(\phi_1,\ldots,\phi_M\) be any finite orthonormal list in \(\cK_N\).
There are even gauge- and \(G\)-invariant vectors \(z_{a,h}\) such that
\begin{equation}\label{count:whole-recovery-estimates}
 \norm{D_u(h,1)z_{a,h}}\longrightarrow0,
 \qquad \ip{z_{a,h}}{z_{b,h}}\longrightarrow\mathbf1_{a=b}.
\end{equation}
They concentrate at zero in physical coordinates. No weighted moment
condition on any \(\phi_a\) is required.
\end{lemma}

\begin{proof}
Put \(d=9(N^2-1)\), and let
\[
 (S_h\phi)(x)=h^{d/6}\phi(h^{1/3}x).
\]
Then \(S_h\) is unitary and
\(D_u(h,0)S_h=h^{1/3}S_hD_u(1,0)\). Choose a fixed compact radial cutoff
\(\chi\) equal to one near zero, and set \(z_{a,h}=\chi S_h\phi_a\).
The vectors \(S_h\phi_a\) concentrate at zero by their \(L^2\) integrability:
for every \(\eta>0\),
\[
 \int_{|x|>\eta}|S_h\phi_a|^2
                  =\int_{|y|>h^{1/3}\eta}|\phi_a(y)|^2\longrightarrow0.
\]
It follows that \(\norm{z_{a,h}-S_h\phi_a}\to0\), proving the Gram
matrix assertion. Lemma~\ref{lem:domains} allows the cutoff product in
the graph domain, and its undeformed charge is
\([D_u,\chi]S_h\phi_a\). The derivative of \(\chi\) is bounded and
supported away from zero, so this norm tends to zero.

The massive multiplier has norm at most \(C|x|\). On the fixed support
of \(\chi\), its norm on \(z_{a,h}\) tends to zero as well: if the support
is contained in \(|x|\le B\), then
\[
 \norm{|x|z_{a,h}}^2
 \le\eta^2\norm{\phi_a}^2+
                   B^2\int_{|x|>\eta}|S_h\phi_a|^2.
\]
Take \(h\to\infty\) and then \(\eta\downarrow0\). This proves the
charge estimate without assuming \(|y|\phi_a\in L^2\).
Proposition~\ref{prop:singlet-parity} makes every \(\phi_a\) even and a
rotation singlet; radial cutoff and dilation preserve these symmetries.

At any fixed level of the proper-partition construction, the compact
center supports avoid collisions and hence avoid zero. Its lifts are
therefore asymptotically orthogonal to all the \(z_{a,h}\). There are
only finitely many such pairs. Include these overlaps in the threshold
selection in Lemma~\ref{count:proper-recovery} to obtain the simultaneous
claim there.
\end{proof}

\subsection{Closing the simultaneous induction}

\begin{proof}[Proof of Theorem~\ref{thm:main}]
Assume that \(\dim\cK_n=1\) and Proposition~\ref{prop:exterior-estimate}
hold for \(2\le n<N\). Section~\ref{sec:profiles} first proves the
exterior estimate at size \(N\), using only proper profiles on compact
annuli. It then constructs the full-collapse profile in \(\cK_N\).
Proposition~\ref{prop:massive-tightness} and
Lemma~\ref{lem:large-coupling-gap} now show that, for all sufficiently
large \(h\), the massive kernel on gauge and \(G\) invariants has
dimension exactly \(p(N)\), with a uniform gap on its complement.

Choose any sequence \(h_j\to\infty\) in this range and an orthonormal
\(p(N)\)-element basis of its massive kernel for each \(j\). Simultaneously
extract profiles from this finite list. Equation~\eqref{prof:polarized-exhaustion}
says that the resulting list is orthonormal in
\begin{equation}\label{count:limit-target}
              \cK_N\ \oplus\!\!\bigoplus_{C\text{ proper}}\C\Phi_C.
\end{equation}
The proper summand has dimension \(p(N)-1\). If \(\cK_N\) were zero,
this space could not contain the orthonormal list. Hence
\(\cK_N\ne\{0\}\).

For the opposite inequality, take an arbitrary finite orthonormal list
\(\phi_1,\ldots,\phi_M\) in \(\cK_N\). Lemmas~\ref{count:proper-recovery}
and \ref{count:whole-recovery} provide
\(L=p(N)-1+M\) asymptotically orthonormal even gauge- and
\(G\)-invariant vectors \(w_{1,h},\ldots,w_{L,h}\) with charge norms
tending to zero. The gap \eqref{count:full-gap} implies
\[
 \norm{w_{a,h}-P_hw_{a,h}}\longrightarrow0
                  \qquad(1\le a\le L).
\]
Their projected Gram matrix therefore tends to the \(L\)-by-\(L\)
identity. For all sufficiently large \(h\) it is positive definite, so
the \(L\) projected vectors are linearly independent in the
\(p(N)\)-dimensional massive kernel. Thus
\[
                       p(N)-1+M\le p(N),\qquad M\le1.
\]
This applies to every finite orthonormal list, and hence proves
\(\dim\cK_N\le1\) without presupposing finite dimensionality.
Together with existence it gives \(\dim\cK_N=1\).

For \(N=2\), every proper internal block has size one. Its kernel factor
is \(\C\), and there are no smaller nonabelian exterior estimates to
assume. The same argument therefore starts the induction. At each
subsequent size the order is: smaller-block kernel and exterior results,
proper-stratum profiles, the current exterior estimate, the full-collapse
profile, and finally the massive dimension comparison. This completes
the simultaneous induction. Lemma~\ref{lem:domains} identifies
\(\cK_N\) with the kernel of the form Hamiltonian in
\eqref{eq:model-form}, proving Theorem~\ref{thm:main} on the entire
configuration space.
\end{proof}
